\scanpage{45}
\mspage{\unsure{19|17|$1'$}}

\uline{Corollaire 3} Soit $\underset{\substack{\text{ouvert}\\\text{0-conn}}}{X\subset X'}$, $X'$ holo $0$-connexe

Alors $X$ est hyperbolique, sauf dans les cas suivants :

a) $X'$ est $\simeq\mathbb P^1_{\mathbb C}$ (resp $\mathbb E^1_{\mathbb C}$, resp. $\mathbb G_{m\,\mathbb C}$), et $\mathrm{Card}\,(X\setminus X')\leq 2$ (resp. $\leq 1$, resp. $=0$)

b) $X'$ est compacte de genre 1, et $X=X'$.

\uline{Remarque} Soit \deleted{\unsure{un|a}} $X$ une surface conforme $0$-connexe (pas néc. orientable). On dit encore que $X$ est hyperbolique si $\widetilde X\simeq\mathbb H$, ou ce qui revient au $=$, si tte \unsure{applic|applic.} conforme $\mathbb C\longrightarrow X$ est constante.

% slanted
\margin{ou encore si le rev. \deleted{\unsure{absol.|orienté|d'orient.}} \inserted{d'orientation} de $X$ (qui est une surface holomorphe,) est hyperbolique.}

Ainsi, les surfaces conformes \inserted{ss bord} $0$-connexes \uline{non} hyperboliques \unsure{sont|2 et}, en plus des quatre cas énumérés dans (i) ci-dessus, les quotients de telles surfaces par des antiinvolutions sans pts fixes. On trouve donc en plus

% archive sheet number: 47
